% Ad Familiares 3.13 (ad-familiares-03-13)
% Source: https://raw.githubusercontent.com/PerseusDL/canonical-latinLit/master/data/phi0474/phi056/phi0474.phi056.perseus-lat1.xml
% Fetched by scripts/fetch_latin.py

% Dateline (Perseus): Scr. Rhodi circ. iv Id. Sext. a. 704 (50).

\ciceroLetterOpener{CICERO APPIO PVLCHRO S.}

\ciceroSection{1} quasi divinarem tali in officio fore mihi aliquando ex petendum studium tuum, sic, cum de tuis rebus gestis agebatur, inserviebam honori tuo; dicam tamen vere: plus quam acceperas reddidisti. quis enim ad me non perscripsit te non solum auctoritate orationis, sententia tua, quibus ego a tali viro contentus eram, sed etiam opera, consilio, domum veniendo, conveniendis meis nullum onus offici cuiquam reliquum fecisse? haec mihi ampliora multo sunt quam illa ipsa, propter quae haec laborantur.

\ciceroSection{2} insignia enim virtutis multi etiam sine virtute adsecuti sunt, talium virorum tanta studia adsequi sola virtus potest. itaque mihi propono fructum amicitiae nostrae ipsam amicitiam, qua nihil est uberius, praesertim in iis studiis, quibus uterque nostrum devinctus est. nam tibi me profiteor et in re publica socium, de qua idem sentimus, et in cotidiana vita coniunctum, quam his artibus studiisque colimus. vellem ita fortuna tulisset, ut, quanti ego omnis tuos facio, tanti tu meos facere posses, quod tamen ipsum nescio qua permotus animi divinatione non despero. sed hoc nihil ad te; nostrum est onus. illud velim sic habeas, quod intelleges hac re novata additum potius aliquid ad meum erga te studium, quo nihil videbatur addi posse, quam quicquam esse detractum Cum haec scribebam, censorem iam te esse sperabam. eo brevior est epistula et ut adversus magistrum morum modestior.
